\documentclass[10pt,a4paper]{amsart}
\usepackage[margin=19mm]{geometry}
\usepackage{amsmath,amssymb,mathtools}
\usepackage[scr=rsfs,cal=euler]{mathalfa}
\usepackage{xfrac}
\usepackage[unicode,hidelinks,bookmarks=false]{hyperref}
\newcommand{\ie}{i.e.\ }
\newcommand{\cf}{cf.\ }
\newcommand{\resp}{respectively\ }
\newcommand{\Subsection}{\subsection}
\providecommand{\nobreakdash}{\nobreak\hbox{-}}
\setlength{\emergencystretch}{2em}
\multlinegap0pt
\usepackage{amssymb}
\newtheorem{lemma}{Lemma}
\title{Reciprocity formulas for certain generalized Hardy-Berndt sums}
\author{Yuan He}
\date{Source: arXiv:2401.00774v6 (CC BY 4.0)}
\begin{document}
\maketitle

\noindent\emph{Source and licence.} Reciprocity formulas for certain generalized Hardy-Berndt sums. Version arXiv:2401.00774v6 (CC BY 4.0), distributed by arXiv under the Creative Commons Attribution 4.0 International licence, http://creativecommons.org/licenses/by/4.0/, as recorded in the arXiv licence metadata for this exact version and shown on the abstract page https://arxiv.org/abs/2401.00774v6. Full licence notice: \texttt{licenses/arxiv-2401.00774-CC-BY-4.0.txt}.\par\smallskip

\noindent\textit{Editorial scope.} Counted unit: the article's lemma lem3.6 (source0.tex lines 654--659). The article's complete original proof of it is delivered with it (source0.tex lines 661--693, one \texttt{proof} environment of 33 lines). No mathematics is written, repaired or re-derived: the statement and the proof are the article's own lines, in the article's own notation, and no retained line is substituted or re-flowed.  Retained uncounted as the source-local context the proof needs: lines 127--132; lines 148--150; lines 567--572; lines 590--592; lines 612--617; lines 640--648; lines 641--643.  Nothing was omitted from any retained span, so the package carries no omission record.  No retained line contains a citation, so the wrapper carries no bibliography and no bibliography entry.  No retained line contains a figure environment, an image-inclusion command or drawing code, so no image is delivered and the package declares no asset dependency.  Editorial formatting only, in addition: an AMS \texttt{amsart} preamble that restates the article's own declarations and macros used by the retained lines, namely the environments \texttt{lemma} on separate counters, together with the standard packages those lines need.  The retained environments print in this document as Lemma 1 in order of appearance, whereas the article prints the same items with the numbering of its own full document.  The complete native source of the article as distributed in the arXiv e-print for this exact version is bundled unchanged as \texttt{sources/152328/source0.tex}.
\begin{equation}\label{eq1.13}
\overline{E}_{0}(x)=\begin{cases}
(-1)^{[x]},  &\text{if $x\in\mathbb{R}\setminus\mathbb{Z}$},\\
0,  &\text{if $x\in\mathbb{Z}$},
\end{cases} \quad\overline{E}_{n}(x)=(-1)^{[x]}E_{n}(\{x\})\quad(n\geq1),
\end{equation}

\begin{equation}\label{eq2.2}
\overline{E}_{n}(x)=\frac{2n!}{(2\pi\mathrm{i})^{n+1}}\sum_{k=-\infty}^{+\infty}\frac{e^{2\pi\mathrm{i}(k-\frac{1}{2})x}}{\bigl(k-\frac{1}{2}\bigl)^{n+1}}.
\end{equation}

\begin{lemma}\label{lem3.4} Let $j\in\mathbb{N}$, $b,r\in\mathbb{Z}$ with $b\not=0$. Then
\begin{equation}\label{eq3.4}
\sum_{d=-\infty}^{+\infty}\frac{1}{\bigl(d+\frac{r}{b}-\frac{1}{2b}\bigl)^{j}}
=\frac{(-1)^{j}(2\pi\mathrm{i})^{j}\mathrm{sgn}(b)}{2(j-1)!b^{1-j}}\sum_{l=1}^{|b|}e^{\frac{2\pi \mathrm{i}lr}{b}-\frac{\pi \mathrm{i}l}{b}}\overline{E}_{j-1}\biggl(\frac{l}{b}\biggl).
\end{equation}
\end{lemma}

\begin{equation}\label{eq3.7}
\eta(-j,x)=\frac{1}{2}E_{j}(x),
\end{equation}

\begin{equation}\label{eq3.11}
\overline{E}_{n}(-x)=\begin{cases}
\overline{E}_{n}(x),  &\text{if $x\in\mathbb{Z}$},\\
(-1)^{n+1}\overline{E}_{n}(x),  &\text{if $x\in\mathbb{R}\setminus\mathbb{Z}$}.
\end{cases}
\end{equation}

\begin{lemma}\label{lem3.5}(Lerch's functional equation) If $0<x<1$ and $0<a\leq1$ then for $s\in\mathbb{C}$,
\begin{equation}\label{eq3.15}
\phi(x,a,1-s)=\frac{\Gamma(s)}{(2\pi)^{s}}\bigl(e^{\frac{\pi \mathrm{i}s}{2}-2\pi\mathrm{i}ax}\phi(-a,x,s)+e^{-\frac{\pi \mathrm{i}s}{2}+2\pi\mathrm{i}a(1-x)}\phi(a,1-x,s)\bigl),
\end{equation}
where $\phi(x,a,s)$ is the Lerch-zeta function given for $x,a\in\mathbb{R}$, $s\in\mathbb{C}$ with $0<a\leq1$ by
\begin{equation*}
\phi(x,a,s)=\sum_{n=0}^{\infty}\frac{e^{2\pi \mathrm{i}nx}}{(n+a)^{s}}\quad(\Re(s)>1).
\end{equation*}
\end{lemma}

\begin{equation}
\phi(x,a,1-s)=\frac{\Gamma(s)}{(2\pi)^{s}}\bigl(e^{\frac{\pi \mathrm{i}s}{2}-2\pi\mathrm{i}ax}\phi(-a,x,s)+e^{-\frac{\pi \mathrm{i}s}{2}+2\pi\mathrm{i}a(1-x)}\phi(a,1-x,s)\bigl),
\end{equation}

\begin{lemma}\label{lem3.6} Let $j\in\mathbb{N}$, $b,r\in\mathbb{Z}$ with $b\not=0$. Then, for $x\in\mathbb{R}$,
\begin{eqnarray}\label{eq3.16}
&&\sum_{d=-\infty}^{+\infty}\frac{e^{2\pi \mathrm{i}dx}}{\bigl(d+\frac{r}{b}-\frac{1}{2b}\bigl)^{j}}\nonumber\\
&&=\frac{(-1)^{j}(2\pi\mathrm{i})^{j}\mathrm{sgn}(b)}{2(j-1)!b^{1-j}}\sum_{l=1}^{|b|}e^{\frac{2\pi \mathrm{i}(l-x)r}{b}-\frac{\pi \mathrm{i}(l-x)}{b}}\overline{E}_{j-1}\biggl(\frac{l-x}{b}\biggl).
\end{eqnarray}
\end{lemma}

\begin{proof}
We first consider the case $x\in\mathbb{Z}$.
Since $l-x$ runs over a complete residue system modulo $|b|$ as
$l$ does. Hence, we know from \eqref{eq2.2} and Lemma \ref{lem3.4} that Lemma \ref{lem3.6} is true in the case when $x\in\mathbb{Z}$. We next discuss the case $x\not\in\mathbb{Z}$. It suffices to prove that Lemma \ref{lem3.6} holds in the case when $0<x<1$. Obviously, in this case, we have
\begin{eqnarray}\label{eq3.17}
&&\sum_{d=-\infty}^{+\infty}\frac{e^{2\pi \mathrm{i}dx}}{\bigl(d+\frac{r}{b}-\frac{1}{2b}\bigl)^{j}}\nonumber\\
&&=e^{-2\pi \mathrm{i}[\frac{r}{b}-\frac{1}{2b}]x}\biggl(\sum_{d=0}^{\infty}\frac{e^{2\pi \mathrm{i}dx}}{(d+\{\frac{r}{b}-\frac{1}{2b}\})^{j}}+(-1)^{j}e^{-2\pi \mathrm{i}x}\sum_{d=0}^{\infty}\frac{e^{-2\pi \mathrm{i}dx}}{(d+1-\{\frac{r}{b}-\frac{1}{2b}\})^{j}}\biggl)\nonumber\\
&&=e^{-2\pi \mathrm{i}[\frac{r}{b}-\frac{1}{2b}]x}\nonumber\\
&&\quad\times\biggl(\phi\biggl(x,\biggl\{\frac{r}{b}-\frac{1}{2b}\biggl\},j\biggl)
+(-1)^{j}e^{-2\pi \mathrm{i}x}\phi\biggl(-x,1-\biggl\{\frac{r}{b}-\frac{1}{2b}\biggl\},j\biggl)\biggl),
\end{eqnarray}
where $\phi(x,a,s)$ is as in \eqref{eq3.15}.
If we take $x=1-\{\frac{r}{b}-\frac{1}{2b}\}$ and $s=j$ and replace $a$ by $x$ in Lemma \ref{lem3.5}, then we have
\begin{eqnarray}\label{eq3.18}
&&\phi\biggl(1-\biggl\{\frac{r}{b}-\frac{1}{2b}\biggl\},x,1-j\biggl)\nonumber\\
&&=\frac{(j-1)!}{(2\pi)^{j}}\biggl(e^{\frac{\pi \mathrm{i}j}{2}-2\pi\mathrm{i}x(1-\{\frac{r}{b}-\frac{1}{2b}\})}\phi\biggl(-x,1-\biggl\{\frac{r}{b}-\frac{1}{2b}\biggl\},j\biggl)\nonumber\\
&&\quad+e^{-\frac{\pi \mathrm{i}j}{2}+2\pi\mathrm{i}x\{\frac{r}{b}-\frac{1}{2b}\}}\phi\biggl(x,\biggl\{\frac{r}{b}-\frac{1}{2b}\biggl\},j\biggl)\biggl).
\end{eqnarray}
Multiplying both sides of \eqref{eq3.18} by $e^{\frac{\pi \mathrm{i}j}{2}-2\pi\mathrm{i}\{\frac{r}{b}-\frac{1}{2b}\}x}$, and it follows from \eqref{eq3.17} that
\begin{equation}\label{eq3.19}
\sum_{d=-\infty}^{+\infty}\frac{e^{2\pi \mathrm{i}dx}}{\bigl(d+\frac{r}{b}-\frac{1}{2b}\bigl)^{j}}
=e^{-\frac{2\pi \mathrm{i}rx}{b}+\frac{\pi \mathrm{i}x}{b}}\frac{(2\pi \mathrm{i})^{j}\phi(1-\{\frac{r}{b}-\frac{1}{2b}\},x,1-j)}{(j-1)!}.
\end{equation}
Note that from the division algorithm, \eqref{eq1.13}, \eqref{eq3.7} and \eqref{eq3.11} we have
\begin{eqnarray}\label{eq3.20}
&&\phi\biggl(1-\biggl\{\frac{r}{b}-\frac{1}{2b}\biggl\},x,1-j\biggl)\nonumber\\
&&=\sum_{l=0}^{|b|-1}\sum_{m=0}^{\infty}\frac{e^{-\frac{2\pi \mathrm{i}(m|b|+l)r}{b}+\frac{\pi \mathrm{i}(m|b|+l)}{b}}}{(m|b|+l+x)^{1-j}}\nonumber\\
&&=\frac{1}{2|b|^{1-j}}\sum_{l=0}^{|b|-1}e^{-\frac{2\pi \mathrm{i}lr}{b}+\frac{\pi \mathrm{i}l}{b}}E_{j-1}\biggl(\frac{l+x}{|b|}\biggl)\nonumber\\
&&=\frac{\mathrm{sgn}(b)}{2b^{1-j}}\sum_{l=0}^{|b|-1}e^{-\frac{2\pi \mathrm{i}lr}{b}+\frac{\pi \mathrm{i}l}{b}}\overline{E}_{j-1}\biggl(\frac{l+x}{b}\biggl)\nonumber\\
&&=\frac{(-1)^{j}\mathrm{sgn}(b)}{2b^{1-j}}\sum_{l=1}^{|b|}e^{\frac{2\pi \mathrm{i}lr}{b}-\frac{\pi \mathrm{i}l}{b}}\overline{E}_{j-1}\biggl(\frac{l-x}{b}\biggl).
\end{eqnarray}
Therefore, inserting \eqref{eq3.20} into \eqref{eq3.19}, we say that Lemma \ref{lem3.6} holds in the case when $x\not\in\mathbb{Z}$. This completes the proof of Lemma \ref{lem3.6}.
\end{proof}

\end{document}
